% ============================================================
% РАЗДЕЛ «ОСНОВНЫЕ ФУНКЦИИ»
% ============================================================

{\color{unidarkgreen}\section[Основные функции]{Основные функции}\label{sec:funcs}}

{\color{unidarkgreen}\subsection{Общие сведения}}

\begin{enumerate}
\item Устройство постоянно отслеживает состояние подключенных к нему аналоговых и дискретных сигналов, а также сигналов портов связи.
\item Устройство периодически производит замер мгновенных значений токов и напряжений при помощи АЦП. Моменты замера АЦП в составе разных каналов синхронизированы, что позволяет исключить погрешность в фазовом сдвиге между отсчетами разных каналов.
\item Из значений, полученных АЦП, при помощи цифровой фильтрации и расчетов выделяются величины (тока, напряжения, фаз, сопротивлений, гармонических составляющих и т.д.), необходимые для использования в алгоритмах устройства. Количество и наименование получаемых величин обусловлены их потребностью в конкретном типоисполнении устройства.
\item Величина аналогового сигнала подводится на измерительный орган функционального блока (ФБ), где сравнивается с уставкой ИО. При достижении величиной значения уставки происходит срабатывание ИО. В ИО всех ФБ присутствует гистерезис для исключения дребезга.
\item \textls[-30]{Сигнал срабатывания ИО, при соблюдении всех необходимых условий, запускает таймеры задержки (при их наличии) срабатывания ФБ. В случае возврата измерительного органа происходит сброс выдержки времени.}
\item После выдержки заданного времени ФБ формируется сигнал срабатывания ФБ.
\item Данный сигнал может быть использован для конфигурирования на выходное реле устройства, передачи в исходящем GOOSE сообщении, фиксации в журнале событий и встроенном регистраторе аварийных событий. Так же данный сигнал может использоваться в качестве входного сигнала других ФБ функциональной схемы устройства (ФСУ).
\item Общий принцип работы устройства и структурная схема работы устройства представлены в \ManualUnitGeneral.
\item \textls[-30]{Оперативное управление ФБ осуществляется посредством виртуальных ключей и виртуальных кнопок.}
\item Оперативный вывод всех функций в составе устройства из работы происходит при появлении активного сигнала «Вывод терминала».
\item Коэффициенты возврата максимальных измерительных органов заданы равными 0,95, коэффициенты возврата минимальных измерительных органов заданы равными 1,05 (если в описании функции не указано другое значение).
\end{enumerate}


\begin{enumerate}[resume]
\item Основные характеристики функциональных блоков:
\begin{itemize}
    \item погрешность измерительных органов тока и напряжения (в том числе ИО по симметричным составляющим) не более 2,5\%, на границе частот 45 Гц и 55 Гц не более 5\%;
    \item время срабатывания ИО тока при повышении тока скачком от 0 до 2Iуст. не более 40 мс;
    \item время возврата реле ИО тока при снижении тока скачком от 2Iуст. до 0 не более 40 мс;
    \item \textls[-15]{время срабатывания ИО напряжения при повышении напряжения скачком от 0 до 2Uуст не более 40 мс;}
    \item \textls[-30]{время возврата ИО напряжения при сбросе напряжения скачком от 2Uуст до 0 должно быть не более 40 мс;}
    \item погрешность выдержки времени -- не более 40 мс при выдержках менее 1,5 с и не более 2\% от уставки при выдержках времени выше 1,5 с.
\end{itemize}
Специфические требования к характеристикам функциональных блоков приведены в описаниях соответствующих функций в данном разделе.
\item Общая схема взаимосвязей между описанными в данном разделе функциями, входными и выходными сигналами устройства приведена в приложении \ref{app:fsu}. Сводный перечень сигналов, формируемых устройством в процессе работы, приведен в таблице \ref{appA:tbl1} приложения \ref{app:signals}.
\end{enumerate}